provides a finite proof-token descent, or enters the typed obligation
normalizer through the one-shot entry component.
\end{proposition}

\begin{proof}
We split according to the phase at $x$.

First suppose $e=1-\alpha(x)$, the $B$-selecting phase. The first source
boundary numerator is divisible by $4$, and
\[
  9J(x)+1-2e=3(3J(x)+1-2e)-2(1-2e)\equiv2\pmod4.
\]
Thus $v=1$ in (1). Writing $u=A(x)$ and using the parity defining the
$B$-selecting phase gives the exact endpoint
\[
  b=3A(x)+1.
\]
The two children of $F$ are
\[
  T=Q_1^g(J(b)),\qquad C=R(T),
\]
and the signed-prefix identity identifies $C$ as the ordinary child of $b$
selected by phase $g$:
\[
  C\in\{A(b),B(b)\}.
\]
At least one of $T,C$ is \DRAW. If $T$ is \DRAW, the first alternative of the
obligation $O(b,g)$ holds and the typed normalizer is initialized. If $C$ is
\DRAW, the other ordinary child of the \WIN\ state $b$ is a genuine \LOSS\
token. For nonexceptional $b$, the ordinary side diamond gives a lower boundary
\WIN\ token. For the exceptional $A$-child orientation, the exceptional side
formula gives an adjacent \DRAW\ frame over the actual \LOSS\ token. In the
sole exceptional $B$-child orientation, the estimate
\[
  C=B(b)\le \frac{A(b)}8\le \frac{3b+1}{16}
\]
combined with $b\le(9x+5)/2$ yields
\[
  \src(C)<\frac{C}{6}\le\frac{3b+1}{96}
  \le\frac{27x+17}{192}<x,
\]
so the retained source strictly decreases.

Now suppose $e=\alpha(x)$, the $A$-selecting phase. Lemma~\ref{lem:source-boundary}
gives
\[
  3J(x)+1-2e=2J(A(x)),
\]
and (1) therefore has $v\ge2$. If $v\ge4$, then
\[
  16J(b)\le 9J(x)+1\le 27x+19.
\]
Since $J(b)\ge3b+1$, we get
\[
  48b+16\le27x+19,
  \qquad
  b\le\frac{9x+1}{16}<x.
\]
Thus any \DRAW\ child of $F$ lies below the retained source.

If $v=2$, the two children of $F$ are exactly
\[
  Q_2^g(J(b)),\qquad Q_1^g(J(b)),
\]
so the obligation $O(b,g)$ holds. If $v=3$, the children form the frame
\[
  Q_3^g(J(b)),\qquad Q_2^g(J(b)).
\]
Reducing (1) modulo $3$ gives
\[
  J(b)\equiv1+g\pmod3,
\]
which is precisely the hidden-parent congruence required for the typed
three/two frame. That frame is therefore normalized by the same source/token
routing used for typed factor frames.
\end{proof}

\section{Typed fixed-fibre normalization}

The local source, factor, and proof-token diamonds used after a typed entry form
a finite routing system. Its control states include $A$-selecting obligations,
$B$-selecting transfers, finite factor forks, high returns, and marked terminal
rows. The important distinction is between retained data and routing data.

The retained data are:
\[
  (\text{outer source},\; \Phi(M),\;\eta),
\]
where $M$ is the finite multiset of carried proof-token heights and
$\eta\in\{1,0\}$ is ordered by $1>0$. The value $\eta=1$ means that the current
outer source/token fibre has not yet initialized the typed normalizer; the
value $\eta=0$ means that the typed normalizer is active.

Routing data include the current source cursor, phase, factor exponent, tail
exponent, and local control state. These quantities may increase numerically.
They do not replace the retained source unless a cited local lemma proves a
strict source comparison.

\begin{proposition}[Typed normalizer]\label{prop:typed-normalizer}
For every typed entry supplied with a retained arithmetic source anchor and the
finite proof-token provenance required by the local diamonds, the equal-retained
data routing relation is well-founded.
\end{proposition}

\begin{proof}[Proof sketch]
The finite routing system has only the following non-strict transitions inside a
fixed retained fibre: bounded canonical $A$-streaks, valuation-two $B$-transfers
with an installed witness, and factor-tail recurrences that carry an existing
finite token. The canonical $A$-streak exits after at most four side tests, since
four consecutive \WIN\ states cannot occur on one $A$-ray. The valuation-two
recycle replaces its witness by a certified lower descendant whenever it repeats.
Marked factor tails either lower their tail counter or replace an existing
\WIN/\LOSS\ token by lower token(s). Hence there is no equal-rank cycle in a
fixed retained fibre. Strict source or token transitions are handled outside the
fibre by the lexicographic rank.
\end{proof}

\begin{lemma}[One-shot entry]\label{lem:eta}
A transition that changes $\eta$ from $1$ to $0$ is strict regardless of the
numerical value of the newly initialized inner routing source. Inside an
unchanged outer source/token fibre, $\eta$ never changes from $0$ back to $1$.
A later raw reset is allowed only after a strict earlier source decrease or a
strict proof-token replacement.
\end{lemma}

\begin{proof}
Use the lexicographic order
\[
  (\text{outer source},\Phi(M),\eta,\text{inner typed rank}).
\]
The component $\eta:1\to0$ strictly decreases after the retained source and
outer token components have been fixed. Because $\eta$ is placed before all
inner routing data, the new inner source may be larger than the old outer source
without being counted as a source update. A change $0\to1$ in the same outer
fibre would increase the rank, so it is not a permitted routing transition.
Only a strict decrease in an earlier component may reset later components.
\end{proof}

\section{Proof of the main theorem}

Assume for contradiction that a \DRAW\ exists and choose the least coefficient
source $s$ of a \DRAW. The source-zero case is handled by the explicit
zero-source constant-tail normalization: a zero-source \DRAW\ reaches an
exponent-one/two boundary with either a strict source exit or a typed finite
terminal token.

Let $s>0$. Take an actual \DRAW\ whose coefficient source is $s$. Long-tail
recurrence reduces it to tail exponent one or two after finitely many
outcome-compatible moves. If a factorful exponent-one state
\[
  Q_1^e(3^kJ(s)),\qquad k>0,
\]
is reached, Proposition~\ref{prop:factorful-entry} gives either finite
provenance for a one-shot typed entry or the factor-free exponent-two state
\[
  Q_2^e(J(s))\text{ \DRAW}.
\]
The latter is handled by Proposition~\ref{prop:base-entry}.

Thus every marked macrostep from a \DRAW\ configuration has one of the following
forms:
\begin{enumerate}[label=(\roman*)]
\item the retained coefficient source strictly decreases;
\item at fixed retained source, an existing finite proof token is replaced by
      certified lower token(s), so Lemma~\ref{lem:token-rank} strictly lowers
      $\Phi$;
\item the retained source and token multiset are unchanged, and the one-shot
      component changes $\eta:1\to0$;
\item $\eta=0$ and the path remains inside the well-founded typed fibre of
      Proposition~\ref{prop:typed-normalizer}.
\end{enumerate}
The lexicographic rank
\[
  (\text{retained source},\Phi(M),\eta,\text{typed fibre rank})
\]
is therefore well-founded.

On the other hand, a \DRAW\ state has no \LOSS\ child and at least one \DRAW\
child. Hence after each finite macrostep one may choose an outcome-compatible
continuing \DRAW\ unless a strict smaller-source contradiction has already
occurred. Repeating this choice would produce an infinite descending path in the
well-founded marked relation, impossible. Hence the conjugated game has no
\DRAW\ positions.

The conjugacy of Section~2 sends the terminal conjugated state $0$ to the
original odd state $1$. Therefore every positive odd starting position has finite
remoteness and optimal play reaches $1$.

\section{Verification boundary and reproducibility}

The proof above is a mathematical proof. The accompanying computational files
are supporting checks with a narrower role:
\begin{itemize}[leftmargin=*]
\item finite Python regressions check the arithmetic identities, valuation rows,
      source inequalities, and finite outcome certificates on stated finite
      ranges;
\item the symbolic global-routing certificate checks the declared typed
      control/rank assembly, but it is not an end-to-end formal proof of the
      theorem;
\item the Lean files in the repository kernel-check selected definitions and
      general well-founded metatheory, not the complete concrete no-\DRAW\
      theorem.
\end{itemize}
A reproducible check should record the repository revision, run the standard
Python check script, run the accompanying arithmetic regression, and read the formal
coverage file before interpreting the machine outputs.

\begin{thebibliography}{9}

\bibitem{Guy1986}
R. K. Guy,
John Isbell's Game of Beanstalk and John Conway's Game of Beans-Don't-Talk,
\emph{Mathematics Magazine} 59(5) (1986), 259--269.

\bibitem{Guy1996}
R. K. Guy,
Unsolved Problems in Combinatorial Games, Problem 42,
in \emph{Games of No Chance}, MSRI Publications 29, Cambridge University Press,
1996.

\bibitem{Althofer2024}
I. Alth\"ofer, M. Hartisch, T. Zipproth,
Analysis of a Collatz Game and Other Variants of the $3n+1$ Problem,
\emph{Advances in Computer Games} (2024), 123--132.

\bibitem{AlthoferPage}
I. Alth\"ofer,
Collatz Problems: The Two-Player $3n+-1$ Game,
\url{https://althofer.de/collatz-prizes.html}.

\bibitem{Repo}
G. Pochuev,
Optimal $3n\pm1$ Game: Proof and Verification Repository,
\url{https://github.com/Grisha-Pochuev/3n-plus-minus-1-game}.

\end{thebibliography}

\end{document}
